%!TeX root=MemoriaTFG.tex

\chapter{Conclusiones}
\label{cap:conclusiones}

Este trabajo ha caracterizado de forma empírica el rendimiento y la
generalización de los principales modelos fundacionales para imagen
dermatológica ---con PanDerm (encoder visual) y DermLIP
(vision-language) como referencias abiertas--- sobre un conjunto común
de datasets externos y de tareas, armonizados con una ontología
jerárquica de tres niveles validada por la Dra.~R.~Taberner. La lectura
de conjunto no designa un modelo vencedor, sino una tesis matizada: el
rendimiento dermatológico depende del nivel ontológico, del régimen de
datos, de la modalidad de entrada y de la estrategia de adaptación, de
modo que el modelo y el protocolo adecuados se eligen según el objetivo
clínico (sección~\ref{sec:disc-especializacion}) y la palanca de mejora
se desplaza del encoder hacia los datos: su calidad, su diversidad y su
cobertura.

Se han alcanzado todos los objetivos propuestos:

\begin{enumerate}
\item \textbf{Reproducir PanDerm sobre hardware no estándar.} El
pipeline de PanDerm Base y Large se ha verificado sobre NVIDIA DGX Spark
(chip Grace Hopper GB10, $128$\,GB de memoria unificada, CUDA~13.0,
precisión BF16, arquitectura aarch64), reproduciendo los rangos
numéricos de Yan et~al.~\cite{panderm2025} sobre los datasets
compartidos y documentando la portabilidad a esta arquitectura.

\item \textbf{Evaluar de forma homogénea catorce modelos fundacionales
sobre doce datasets.} El benchmark aplica cuatro protocolos principales
bajo configuración estándar (sección~\ref{sec:protocolos}) a PanDerm,
DINOv2, tres variantes de DermLIP, BiomedCLIP, SigLIP-Large,
SAM2.1-Large, GPT-4o, Gemini~2.5~Pro, MedGemma y BLIP-2, más dos CNN
supervisadas como referencia no fundacional, con desglose por dataset y
métrica (tabla~\ref{tab:lp-panderm}, tabla~\ref{tab:benchmark}) y la
comparación frente a los LLMs multimodales en la
sección~\ref{sec:llm-results} (tabla~\ref{tab:llm-comparativa}).

\item \textbf{Construir una ontología jerárquica L1/L2/L3.} Cuatro
categorías etiológicas, $43$ subcategorías clínicas y $367$ diagnósticos
codificados en SNOMED~CT y CIE-10, que armonizan once de los doce
datasets ($72\,654$ imágenes sobre vocabulario común) y habilitan
comparaciones inter-dataset defendibles en L1 y L2.

\item \textbf{Entregar el dataset DermapixelAI~1.0.} Primer dataset
dermatológico verificado en español con imagen, metadatos y texto
narrativo por caso ($1\,089$ imágenes, $669$ casos, mediana de $223$
palabras), particionado case-aware ($891 / 157 / 41$) y distribuido bajo
licencia CC-BY-NC-SA~4.0 con datasheet formal (repositorio reproducible,
\ref{cap:repo}, secciones~\ref{sec:repo-ablaciones}
y~\ref{sec:repo-datasheet}).

\item \textbf{Auditar el solapamiento con el preentrenamiento.} Un
procedimiento por hash MD5 sobre los doce datasets identifica Dermnet
como único conjunto con solapamiento exacto ($100\%$ contenido en
Derm1M) y declara el caveat correspondiente sobre las cifras zero-shot
de la familia DermLIP en ese dataset.

\item \textbf{Dejar un marco metodológico reproducible.} La articulación
de todo lo anterior bajo un protocolo común, con código, splits
canónicos y outputs documentados (repositorio reproducible,
\ref{cap:repo}, sección~\ref{sec:repo-reproducibilidad}), es en sí misma
una contribución transferible a futuros modelos fundacionales
dermatológicos.
\end{enumerate}

\label{sec:contrib-anexos}%
A estas contribuciones del cuerpo se suman varias complementarias,
detalladas en los anexos: un análisis de equidad por fototipo de cinco
encoders sobre Fitzpatrick17k completo (tabla~\ref{tab:equidad-fototipo});
la caracterización de la sensibilidad al tokenizer en DermLIP ---hasta
$48{,}8$\,pp de AUROC entre la peor y la mejor formulación zero-shot---;
la adaptación de SAM2.1 al dominio dermatoscópico por fine-tuning del
decoder (Dice $0{,}947$ sobre ISIC2018); un ensemble de seguridad para
melanoma que alcanza recall del $100\%$ ($70$ de $70$, top-3) sobre el
test de HAM10000; la interpretabilidad por Sparse Autoencoders validada
frente a SkinCon~\cite{skincon} (AUROC media $0{,}880$); y, sobre todo,
\textbf{Dermapixel R0}, la adaptación al castellano de PanDerm Large con
una rama supervisada (head L2 con LoRA, BAcc $0{,}363 \pm 0{,}007$,
sección~\ref{sec:dermapixel-spanderm-v0}) y una rama contrastiva
multimodal explorada como metodología (sección~\ref{sec:spanderm-clip}).

\section{Hallazgos metodológicos transferibles}

Más allá del benchmark concreto, el trabajo deja varias lecciones útiles
para quien quiera montar un sistema parecido. La principal ya se ha
contado: no hay un método mejor para todo, así que la decisión correcta
es elegir el componente según la tarea. De ahí se siguen reglas
prácticas. Con pocos datos etiquetados (del orden de unos cientos de
imágenes) conviene quedarse con el encoder congelado y un linear probe
encima, porque el fine-tuning denso tiende a sobreajustar; con muchos
datos, el fine-tuning sí compensa. En el régimen zero-shot, lo decisivo
no es la imagen sino que el texto y la imagen «hablen el mismo idioma»:
conviene comprobar antes, sobre datos del dominio, que el tokenizer
encaja con el encoder. En segmentación, basta con ajustar el decoder de
un modelo promptable grande dejando el encoder visual congelado. Y en la
long tail ---esos diagnósticos con poquísimos ejemplos--- la recuperación
por similitud (k-NN sobre embeddings, imagen contra imagen) rinde mejor
que un clasificador entrenado, que se queda sin muestras. Por último, dos
cautelas de método: la equidad no se puede leer sola ---un modelo malo por
igual en todos los fototipos parece equitativo sin serlo---, y la
interpretabilidad por Sparse Autoencoders muestra que el encoder ya usa
internamente conceptos del dermatólogo, lo que abre una vía de
trazabilidad clínica.

\section{Aportaciones a la práctica clínica}

Para un servicio como el del HUSLL, el trabajo deja piezas directamente
aprovechables. La primera es una arquitectura sencilla y defendible:
PanDerm Large congelado con un linear probe cuando hay pocos datos, o
ajustado cuando hay muchos; SAM2.1-Large adaptado cuando haga falta
localizar la lesión; y, para el cribado de melanoma con prioridad de no
dejar pasar ninguno, el ensemble de seguridad que alcanza el máximo
recall. La segunda es la ontología L1/L2/L3, una columna vertebral
semántica compatible con la nomenclatura clínica (SNOMED~CT, CIE-10) que
permite enseñar al médico el diagnóstico fino junto a su agrupación, de
modo que el sistema sea robusto en la long tail sin perder especificidad
cuando la hay. La tercera es el propio dataset DermapixelAI~1.0, un
recurso público en español reutilizable por otros grupos. Y la cuarta es
una buena práctica de honestidad: declarar siempre el posible solapamiento
entre los datos de evaluación y los de preentrenamiento, para que las
cifras no engañen.

\section{Trabajo futuro}
\label{sec:lineas-abiertas}

Este trabajo es una primera etapa; el recorrido deja abiertas cuatro
líneas con base técnica y clínica concreta, que le dan continuidad
natural.

\subsection{Un Dermapixel en castellano entrenado desde cero}
\label{sec:dermapixel-futuro}
Dermapixel R0 es una adaptación ligera (LoRA) sobre un modelo preentrenado
en inglés; demuestra que la especialización al castellano es viable con
recursos modestos, pero hereda las limitaciones de partir de un encoder
ajeno. El salto natural es entrenar \emph{desde cero} un modelo
fundacional dermatológico en castellano, con la misma filosofía de
\emph{masked latent modeling} de PanDerm pero sobre material propio: todo
el archivo del banco hospitalario del HUSLL ---del orden de $22\,000$
estudios con sus historias clínicas asociadas, previa aprobación del
comité ético---, que aporta a la vez imágenes y texto clínico real en
español a una escala muy superior a la de DermapixelAI~1.0. Las historias
clínicas no son un metadato accesorio, sino el eje de texto que permite
aprender la asociación imagen-lenguaje en el propio idioma y dominio, justo
donde la rama contrastiva exploratoria de R0 se quedó corta por falta de
pares (sección~\ref{sec:spanderm-clip}). Ese volumen es el que permitiría
abordar la clasificación L3 directa, hoy limitada por la long tail, y
construir un espacio texto-imagen en castellano con opciones reales de
generalizar. Como el dato clínico no puede salir del centro, el
entrenamiento se plantea sobre la infraestructura propia ---la DGX Spark ya
validada en este trabajo---, con las garantías de privacidad que exige.
La cola de diagnósticos raros se cubriría con anotación experta dirigida
por la Dra.~R.~Taberner mediante aprendizaje activo
\emph{in-domain}\label{sec:linea-active-learning}, priorizando las clases
con menos ejemplos en lugar de añadir datos genéricos ---estrategia que,
como mostró el experimento de escala externa, no rellena la long tail---.

\subsection{Conceptos clínicos propios, Concept Bottleneck y
explicabilidad}
Una de las líneas más prometedoras nace de la interpretabilidad por Sparse
Autoencoders. El trabajo ya mostró que las representaciones de PanDerm
codifican, sin supervisión, los conceptos clásicos de la dermatoscopia; el
siguiente paso lo convierte en herramienta clínica. Con la
Dra.~R.~Taberner ya se ha definido un primer vocabulario de conceptos
dermatoscópicos propios ---un avance en sí mismo---, de los cuales varios
no figuran en las escalas clásicas en inglés
(sección~\ref{sec:sae-derm7pt}): son precisamente los que justifican un
vocabulario en español pensado para la práctica real. El plan es anotar
imágenes sobre ese vocabulario ---con una interfaz sencilla y autocontenida
que no exige instalación al especialista---, entrenar las \emph{features}
del SAE para detectarlos y, a partir de ahí, montar un \emph{Concept
Bottleneck Model}: que el sistema estime primero esos conceptos con nombre
clínico y decida \emph{a través} de ellos, no como caja negra. Aunque esa
transparencia cueste algo de precisión ---como se cuantificó sobre el
\emph{Seven-Point Checklist}---, es la vía para que la salida sea explicable
y auditable por el dermatólogo, e incorpora la explicabilidad de serie en
lugar de añadirla a posteriori. La confianza del clínico, y no solo el
acierto, es lo que decide la adopción de un sistema de apoyo.

\subsection{Validación clínica prospectiva}
\label{sec:linea-prospectivo}
Todo lo anterior se mide hoy sobre archivos curados, no sobre el flujo
asistencial real. El paso ineludible antes de cualquier despliegue es una
validación prospectiva sobre cohorte real del HUSLL, con aprobación del
comité ético de investigación, que mida la \emph{utilidad clínica} ---no
solo el rendimiento sobre datasets--- en el circuito de teledermatología:
imágenes adquiridas con dispositivos heterogéneos, iluminación y encuadre
sin estandarizar, y la decisión del facultativo como referencia. Un diseño
natural es un estudio de lectura que compare el criterio del especialista
con y sin el apoyo del sistema, midiendo no solo exactitud sino tiempo de
lectura, concordancia y derivaciones evitadas. Esta evidencia es, además,
la base imprescindible para una eventual certificación como producto
sanitario.

\subsection{IA agéntica y modelos multimodales en español}
\label{sec:linea-agentica}
A medio plazo, los componentes desarrollados (encoder PanDerm, ontología
L1/L2/L3, recuperación visual imagen$\to$imagen, clasificador L2 en
castellano, conceptos y un LLM multimodal en local) encajan en un sistema
\emph{agéntico}: varios agentes especializados ---uno por tarea--- que
trabajan en paralelo y cuyas salidas se combinan en un \emph{ensemble} bajo
control del facultativo, en lugar de delegar todo en un único modelo.
Frente al uso directo de un LLM generalista (sección~\ref{sec:disc-llm}),
reparte la decisión en piezas verificables y reduce el sesgo léxico. La
dirección la marcan agentes clínicos recientes como AMIE, de Google
DeepMind, que alcanzan un rendimiento no inferior al de médicos de atención
primaria en estudios simulados~\cite{deepmind_aicoclinician,amie2026}.
Dentro de esta línea quedan dos piezas: la rama generativa ---modelos tipo
BLIP-2 (sección~\ref{sec:llm-results}), que podrían generar descripciones
clínicas en castellano a partir de la imagen--- y el reentrenamiento de la
rama de texto: dada la sensibilidad al tokenizer
(sección~\ref{sec:disc-zs}), continuar el preentrenamiento del encoder
textual sobre los casos en español ---y, a mayor escala, sobre el banco
HUSLL--- debería mejorar el rendimiento \emph{zero-shot} y la recuperación
en castellano.
